\documentclass[11pt]{article}

\usepackage{amsmath}
\usepackage{amssymb}
\usepackage{geometry}
\usepackage{graphicx}
\usepackage[hidelinks]{hyperref}
\usepackage{pgfplots}
\pgfplotsset{compat=1.18}
\geometry{margin=1in}

\title{Power and Precision in Balanced vs.\ Unbalanced Case--Control Studies}
\author{Dieter Best, Bioinformatics Data Scientist, \\
  dieterbest@cdph.ca.gov \\
\url{https://github.com/ranlib/power_and_precision_calculation_for_case_control_studies}
}
\date{\today}

\begin{document}

\maketitle

\section{Overview}

We compare study designs in terms of:
\begin{itemize}
    \item \textbf{Statistical power}
    \item \textbf{Precision (confidence interval width)}
\end{itemize}

Key message: Increasing controls yields diminishing returns for both power and precision, while increasing cases is far more effective.

\section{Model}

\begin{align*}
p_1 &= \frac{OR \cdot p_0}{1 - p_0 + OR \cdot p_0} \\
\mathrm{Var}(\log OR) &= \frac{1}{n_1 p_1} + \frac{1}{n_1(1-p_1)} + \frac{1}{n_0 p_0} + \frac{1}{n_0(1-p_0)}
\end{align*}

\section{Parameters}

\begin{itemize}
    \item $OR = 2$
    \item $p_0 = 0.2$
    \item $\alpha = 0.05$
    \item $z_{0.975} = 1.96$
\end{itemize}

\newpage

\section{Results}

\subsection{Power}

\begin{figure}[htbp]
    \centering
    \includegraphics[width=0.9\textwidth]{power_plot.png}
    \caption{Power vs. control-to-case ratio}
    \label{fig:power}
\end{figure}

\newpage

\subsection{Confidence Interval Width}

\begin{figure}[htbp]
    \centering
    \includegraphics[width=0.9\textwidth]{ci_width_plot.png}
    \caption{Power vs. control-to-case ratio}
    \label{fig:power}
\end{figure}


\section{Interpretation}

\begin{itemize}
    \item Increasing controls reduces CI width initially, but quickly plateaus
    \item Increasing cases leads to continuous improvement in precision
    \item For fixed total sample size, balanced designs minimize CI width
\end{itemize}

\section{Conclusion}

\begin{itemize}
    \item Power and precision are primarily driven by the number of cases
    \item Additional controls provide diminishing returns
    \item Balanced designs are optimal when feasible
\end{itemize}

\section{When Extra Controls Do Help Meaningfully}

While increasing the number of controls in a case--control study typically yields diminishing returns in statistical power, there are important scenarios in which adding controls can still provide meaningful benefits.

\subsection{Low Baseline Exposure Probability}

When the exposure probability in controls ($p_0$) is very small, the variance contribution from the control group becomes large:
\[
\mathrm{Var}(\log OR) 
= \frac{1}{n_1 p_1} + \frac{1}{n_1 (1 - p_1)} 
+ \frac{1}{n_0 p_0} + \frac{1}{n_0 (1 - p_0)}
\]
In such settings, increasing the number of controls can substantially reduce variance and improve power, even beyond the usual 4:1 control-to-case ratio rule of thumb.

\subsection{Rare Outcomes with Limited Cases}

In many real-world studies (e.g., rare diseases), the number of cases is fixed and difficult to increase. In this situation:
\begin{itemize}
    \item Increasing cases is not feasible
    \item Increasing controls is often inexpensive and practical
\end{itemize}

Although gains in power plateau, additional controls still provide incremental improvements and may be the only viable strategy.

\subsection{Improved Precision of Effect Estimates}

Even when power gains are small, increasing the number of controls reduces the width of the confidence interval:
\[
\text{CI width} \propto \sqrt{\mathrm{Var}(\log OR)}.
\]

This is important when:
\begin{itemize}
    \item Estimation (effect size precision) matters more than hypothesis testing
    \item Results are used for downstream modeling or meta-analysis
\end{itemize}

\subsection{Control Reuse and External Data Sources}

In some designs, controls can be obtained at very low cost:
\begin{itemize}
    \item Public databases
    \item Biobanks
    \item Previously collected cohorts
\end{itemize}

When controls are effectively ``free,'' it can be advantageous to include many of them, even if statistical gains are modest.

\subsection{Adjustment for Confounding}

Larger control groups can improve:
\begin{itemize}
    \item Stability of regression models
    \item Estimation of covariate effects
    \item Matching or stratification procedures
\end{itemize}

This is especially relevant in high-dimensional settings or when adjusting for many confounders.

\subsection{Summary}

Although balanced designs are statistically optimal for a fixed total sample size, increasing the number of controls can still be beneficial when:
\begin{itemize}
    \item Exposure is rare
    \item Cases are limited
    \item Precision (not just power) is important
    \item Controls are inexpensive or readily available
\end{itemize}

Thus, the choice of control-to-case ratio should reflect not only statistical efficiency, but also practical constraints and study objectives.

\appendix
\section{Interpretation of Model Parameters}

This appendix describes the meaning of the parameters used in the power and sample size calculations for the case--control study.

\begin{itemize}

\item \textbf{Odds Ratio ($OR = 2$)} \\ The odds ratio represents the
  strength of association between the exposure and the outcome. An
  odds ratio of 2 indicates that the odds of exposure among cases are
  twice the odds of exposure among controls. Formally,
\[
OR = \frac{ \frac{p_1}{1 - p_1} }{ \frac{p_0}{1 - p_0} },
\]
where $p_1$ is the exposure probability among cases and $p_0$ is the
exposure probability among controls.

\vspace{0.5em}

\item \textbf{Control Exposure Probability ($p_0 = 0.2$)} \\
This is the assumed prevalence of the exposure in the control group,
which serves as a proxy for the exposure prevalence in the underlying
source population. The value of $p_0$ determines the variance of the
estimated odds ratio and therefore directly influences statistical
power.

\vspace{0.5em}

\item \textbf{Significance Level ($\alpha = 0.05$)} \\
The significance level $\alpha$ is the probability of a Type I error,
i.e., rejecting the null hypothesis when it is true. A two-sided
$\alpha = 0.05$ corresponds to a 5\% chance of falsely declaring an
association when none exists.

\vspace{0.5em}

\item \textbf{Critical Value ($z_{0.975} = 1.96$)} \\
The value $z_{0.975}$ is the 97.5th percentile of the standard normal
distribution. It is used in two-sided hypothesis testing and
confidence interval construction with $\alpha = 0.05$, since
\[
z_{1 - \alpha/2} = z_{0.975} \approx 1.96.
\]
This value defines the rejection threshold for the test statistic and
determines the width of the confidence interval.

\end{itemize}

\noindent Together, these parameters determine the expected
distribution of the log odds ratio estimator, its variance, and the
resulting statistical power and confidence interval width in the study
design.

This topic is discussed in \cite{ferdous2022microbiome_power}.

\bibliographystyle{plain}
\bibliography{references}

\end{document}
